\section{Interpretation after the configuration check}
\label{sec:torque-ambiguity-discussion}
Torque observes a normal projection and leaves a radial slope free. Adding
a smooth radial potential can preserve reflection, reciprocity and positive
curvature while leaving every torque value unchanged. Zero discrepancy is
therefore insufficient for full flux reconstruction, even with an ideal sensor.
Pointwise minimum change chooses a representative; it need not retain curl-free
structure across the current plane.

The six-phase interpretation resolves the leading scale disagreement in the
real case. It does not qualify the independent CAN observation as electromagnetic
truth. If its shaft meaning were established, the signed rotor balance would be
\begin{equation}
 \rotorinertia\dot{\mechanicalspeed}
 =M_{\rm em}-\shafttorque-\mechanicallosstorque.
 \label{eq:shaft-torque-balance}
\end{equation}
Even at steady speed mechanical losses separate transmitted and electromagnetic
torque. Equal represented system currents, matched acquisition windows and
terminal-voltage meaning are additional requirements. Sensor offset produces
inverse-current amplification, gain a torque-proportional term, and current
bias changes both the map argument and lever arm. A coherent common rotation
preserves oriented area, while inconsistent frames can change it.

The synthetic ranks apply to their stated operators and support. Real
joint electrical-mechanical estimation needs covariance, parameter scaling
and nuisance columns. A radial magnetic or electrical observation remains
necessary even after the torque channel is fully qualified.
